\section{Сучасна постквантова асиметрична криптографія}

\subsection{Чому поява достатньо потужного квантового комп'ютера створює загрозу для RSA та криптосистем на еліптичних кривих?}

Безпека RSA та криптосистем на еліптичних кривих ґрунтується на математичних задачах, які складно розв'язати за допомогою 
звичайних комп'ютерів. Для RSA такою задачею є факторизація великих цілих чисел, а для криптосистем на еліптичних 
кривих— задача дискретного логарифмування.

Поява достатньо потужного квантового комп'ютера змінює цю ситуацію. Алгоритм Шора дозволяє розв'язувати обидві задачі значно ефективніше, 
ніж класичні алгоритми. Тому квантовий комп'ютер, який матиме достатню кількість стабільних кубітів, зможе отримувати секретні 
ключі на основі відкритих даних. Це поставить під загрозу конфіденційність інформації та надійність цифрових підписів, створених 
за допомогою цих алгоритмів.

\subsection{Пояснити поняття постквантової криптографії (Post-Quantum Cryptography, PQC).}

Постквантова криптографія — це напрям криптографії, який розробляє алгоритми, стійкі як до звичайних, так і до квантових комп'ютерів. 
Її мета — забезпечити захист інформації навіть тоді, коли з'являться достатньо потужні квантові обчислювальні системи.

Важливо, що такі алгоритми не потребують квантового комп'ютера для своєї роботи. Вони виконуються на звичайних комп'ютерах, 
але використовують математичні задачі, для яких наразі не відомо ефективних способів розв'язання навіть за допомогою квантових алгоритмів.

\subsection{Визначити призначення ML-KEM та ML-DSA і пояснити, які традиційні криптографічні задачі вони повинні виконувати.}

ML-KEM (Module-Lattice-Based Key-Encapsulation Mechanism) — це механізм інкапсуляції ключів, призначений для безпечного 
встановлення спільного секретного ключа між двома сторонами. За своїм призначенням він виконує задачу, подібну до традиційних 
механізмів обміну ключами, наприклад Diffie--Hellman.

ML-DSA (Module-Lattice-Based Digital Signature Algorithm) — це алгоритм цифрового підпису, який дозволяє підтвердити 
справжність повідомлення та перевірити, чи не було його змінено. Таким чином, він виконує ті самі основні функції, 
що й традиційні алгоритми цифрового підпису, наприклад RSA або ECDSA.

Обидва алгоритми належать до постквантової криптографії та ґрунтуються на математичних задачах, пов'язаних із ґратками.

\subsection{Пояснити різницю між механізмом встановлення ключа (KEM) та алгоритмом цифрового підпису.}

KEM та цифровий підпис призначені для розв'язання різних криптографічних задач. Механізм встановлення ключа використовується 
для отримання спільного секрету між двома сторонами. Надалі цей секрет можна застосувати для симетричного шифрування даних.

Цифровий підпис, навпаки, не призначений для створення спільного ключа або приховування інформації. Його використовують для 
перевірки цілісності повідомлення та підтвердження того, що воно було підписане власником відповідного приватного ключа.

Отже, KEM забезпечує встановлення секретного ключа, а цифровий підпис — автентичність і цілісність даних.

\subsection{Пояснити принцип гібридного встановлення ключів, при якому одночасно використовуються традиційний та постквантовий механізми.}

Гібридне встановлення ключів передбачає одночасне використання двох механізмів: традиційного, наприклад ECDH, та постквантового, 
наприклад ML-KEM. Кожен із них формує власний спільний секрет, після чого ці секрети об'єднуються за допомогою спеціальної функції виведення ключів.

У результаті отримується спільний ключ, який використовується для подальшого шифрування даних. Такий підхід дозволяє поєднати перевірену 
часом класичну криптографію з новими постквантовими алгоритмами. Якщо один із механізмів виявиться вразливим, інший за умови правильної 
реалізації може забезпечити додатковий рівень захисту.

\subsection{Пояснити, чому перехід інформаційних систем на постквантові алгоритми може відбуватися поступово.}

Перехід на постквантову криптографію не можна виконати миттєво, оскільки сучасні інформаційні системи використовують 
велику кількість різних криптографічних алгоритмів, протоколів і програмних засобів. Їх заміна потребує часу, тестування 
та перевірки сумісності з наявним обладнанням і програмним забезпеченням.

Крім того, постквантові алгоритми можуть мати інші вимоги до розміру ключів, підписів і обсягу переданих даних. 
Тому їх впровадження може вимагати змін у роботі окремих систем. Поступовий перехід дозволяє спочатку використовувати гібридні рішення, 
перевірити їхню надійність і лише після цього повністю перейти на нові алгоритми.
